\documentclass[a4paper, 12pt]{article}
% math symbols
\usepackage{amssymb}
\usepackage{amsmath}
\usepackage{mathrsfs}
\usepackage{mathseries}


\usepackage[margin = 2cm]{geometry}

\tolerance = 1000
\emergencystretch = 0.74cm



\pagestyle{empty}
\parindent = 0mm

\renewcommand{\coursetitle}{DM/ML}
\setcounter{curtask}{1}

\setmathstyle{\includegraphics[scale = 0.05]{../pics/utia-rest.png}}{Мат. логика}{3 курс}


\begin{document}

\libproblem{math-logic}{discrete-lo-categorical}
\libproblem{math-logic}{th-z-le-0-1-dec}
\libproblem{math-logic}{th-eq-arity-1-dec}
\libproblem{math-logic}{th-t-dec-t-dec}

\libdefinition{ariphmetic-predicate}

\libproblem{math-logic}{ariphmetic-predicate}

\libdefinition{associative-calculus}

\task{
    Покажите, что существует такое двустороннее ассоциативное исчисление, для которого вопрос о
    выводимости строки $x$ из строки $y$ является алгоритмически неразрешимым.
}

\dzcomment{
    В данной задаче можно пользоваться невычислимостью одностороннего ассоциативного исчисления.
}

\task{
    Пусть $T$~--- разрешимая непротиворечивая теория. Покажите, что существует полное разрешимое
    непротиворечивое расширение теории $T$.
}

\libproblem{computability}{taut-enumerable}






\end{document}



%%% Local Variables:
%%% mode: latex
%%% TeX-master: t
%%% End:
